% ============================================================================
% Fichier  : partie3/P03_C10_active_directory.tex
% Titre    : Active Directory
% Auteur   : MINKA MI NGUIDJOÏ Thierry Emmanuel
% Version  : 1.0 -- novembre 2026
% ============================================================================
% Sources : RFC 4120 (Kerberos) ; documentation publique de Microsoft (comptes
% de service gérés, modèle d'administration par niveaux, LAPS, groupe
% Protected Users) ; MITRE ATT&CK (T1558.003 Kerberoasting, T1558.004 AS-REP
% Roasting). Opération LOGONE : cas fictif, aucun code d'attaque.
% ============================================================================

\chapter{Active Directory}
\label{chap:ad}
\chaptermark{Active Directory}

\epigraph{%
    « Dans un domaine, l'attaquant ne cherche pas la faille~;\\
    il cherche le chemin. »%
}{Principe du cours}

\niveauChapitre{Expert}{Chapitre~\ref{chap:reseau}~; cours SEC3266,
chapitre sur les identités.}

\begin{boxobjectifs}
\begin{itemize}
    \item Décrire les objets d'un annuaire Active Directory et les groupes à
          privilèges.
    \item Expliquer l'obtention de tickets Kerberos et ses détournements
          (kerberoasting, AS-REP roasting).
    \item Raisonner en chemins d'attaque dans un graphe de droits.
    \item Appliquer le modèle d'administration par niveaux.
    \item Détecter et corriger les faiblesses les plus exploitées.
\end{itemize}
\end{boxobjectifs}

\bigskip

% ============================================================================
\section*{Question directrice : comment un compte d'utilisateur ordinaire devient-il administrateur du domaine en cinq heures~?}
% ============================================================================

Le 28~octobre, Mme~NGO~MBOCK part du compte d'utilisateur ordinaire fourni
par LiZbiZ. Elle n'exploite aucune faille logicielle. Elle interroge
l'annuaire, comme tout utilisateur peut le faire, demande un ticket pour un
service, casse hors ligne le mot de passe d'un compte de service, et
découvre que ce compte est administrateur du domaine. Ce chapitre, sur deux
séances, déroule ce chemin et ses parades.

\medskip
\textbf{Plan du chapitre.}
\begin{enumerate}
    \item L'annuaire et ses privilèges (section~\ref{sec:c10-annuaire}).
    \item Kerberos en bref (section~\ref{sec:c10-kerberos}).
    \item Le kerberoasting (section~\ref{sec:c10-kerberoasting}).
    \item L'AS-REP roasting (section~\ref{sec:c10-asrep}).
    \item Les chemins d'attaque (section~\ref{sec:c10-chemins}).
    \item Le modèle par niveaux et les défenses (section~\ref{sec:c10-defense}).
\end{enumerate}

\separateur

% ============================================================================
\section{L'Annuaire et ses Privilèges}
\label{sec:c10-annuaire}
% ============================================================================

\begin{boxdef}
\termecle{Active Directory}~: service d'annuaire de Microsoft qui gère, dans
un \emph{domaine}, les comptes, les groupes, les ordinateurs et les
stratégies, et qui authentifie les utilisateurs par Kerberos. Les
\termecle{contrôleurs de domaine} en hébergent la base.
\end{boxdef}

\begin{tableaucours}{L{4.4cm}L{7.8cm}}{%
    \tableauHeader{Objet ou groupe} &
    \tableauHeader{Pourquoi il intéresse l'attaquant}
}
Admins du domaine & Contrôle total du domaine \\
Administrateurs de l'entreprise & Contrôle de toute la forêt \\
Opérateurs de sauvegarde, de serveur & Accès étendu, souvent oublié \\
Comptes de service & Mots de passe rarement changés, droits souvent
excessifs \\
Stratégies de groupe & Exécution de code sur toutes les machines liées \\
Listes de contrôle d'accès & Un droit d'écriture sur un objet peut suffire à
en prendre le contrôle \\
\end{tableaucours}
\vspace{-0.8em}
\begin{center}{\small\itshape Tableau~10.1. Ce que vise l'attaquant dans l'annuaire}\end{center}

Tout utilisateur authentifié peut, par défaut, lire l'essentiel de
l'annuaire~: comptes, groupes, attributs. C'est utile au fonctionnement, et
c'est ce qui rend la reconnaissance interne si facile.

% ============================================================================
\section{Kerberos en Bref}
\label{sec:c10-kerberos}
% ============================================================================

\begin{tableaucours}{cL{5.8cm}L{5.2cm}}{%
    \tableauHeader{Étape} &
    \tableauHeader{Échange} &
    \tableauHeader{Secret utilisé}
}
1 & L'utilisateur demande un ticket d'accès général (TGT) au centre de
distribution & Clé dérivée du mot de passe de l'utilisateur
(pré-authentification) \\
2 & Il présente le TGT et demande un ticket de service pour un nom de
service (SPN) & Clé du centre de distribution \\
3 & Le centre renvoie un ticket de service chiffré & Clé dérivée du mot de
passe du \emph{compte de service} \\
4 & L'utilisateur présente le ticket au service & Le service le déchiffre \\
\end{tableaucours}
\vspace{-0.8em}
\begin{center}{\small\itshape Tableau~10.2. Les échanges Kerberos (d'après la RFC~4120)}\end{center}

% ============================================================================
\section{Le Kerberoasting}
\label{sec:c10-kerberoasting}
% ============================================================================

\begin{boxdef}
\termecle{Kerberoasting}~: technique (ATT\&CK T1558.003) par laquelle un
utilisateur authentifié demande des tickets de service pour des comptes
associés à un SPN, puis tente de retrouver hors ligne leur mot de passe à
partir du ticket chiffré.
\end{boxdef}

Demander un ticket de service est un geste normal~: n'importe quel
utilisateur le fait. La faiblesse est ailleurs~: un compte de service dont
le mot de passe est court ou prévisible se casse hors ligne, sans laisser de
trace d'échec.

\begin{boxscenario}{Opération LOGONE. Le compte de la facturation}
\textbf{Constatations} (28~octobre).
\begin{itemize}
    \item 9~h~10~: interrogation de l'annuaire~; 6~comptes ont un SPN.
    \item 9~h~25~: tickets de service demandés pour ces 6~comptes.
    \item 11~h~40~: le mot de passe de « svc\_sql\_facturation » est retrouvé
          hors ligne~: 11~caractères, formé du nom de l'application, d'un
          symbole et d'une année.
    \item 14~h~05~: ce compte est membre des Admins du domaine (déjà relevé
          à la revue des comptes de mars, cours SEC3266).
    \item 14~h~10~: preuve capturée, test arrêté, contacts prévenus.
\end{itemize}

\medskip
\textbf{Analyse} (confiance élevée). Aucune faille logicielle~: trois
faiblesses de configuration se combinent (mot de passe faible, mot de passe
jamais changé, droits excessifs). Le SOC n'a rien vu~: aucune règle ne
surveillait les demandes de tickets.
\end{boxscenario}

% ============================================================================
\section{L'AS-REP Roasting}
\label{sec:c10-asrep}
% ============================================================================

Lorsqu'un compte a la \emph{pré-authentification Kerberos désactivée},
n'importe qui peut demander pour lui une réponse chiffrée avec une clé
dérivée de son mot de passe, et tenter de la casser hors ligne (ATT\&CK
T1558.004). Cette option, parfois activée pour d'anciennes applications,
devrait être exceptionnelle et tracée. Chez LiZbiZ, deux comptes l'avaient
encore, dont un compte de test oublié.

% ============================================================================
\section{Les Chemins d'Attaque}
\label{sec:c10-chemins}
% ============================================================================

L'annuaire forme un \termecle{graphe}~: des utilisateurs membres de groupes,
des groupes ayant des droits sur des objets, des sessions ouvertes sur des
machines. Un chemin d'attaque est une suite d'arêtes qui mène d'un compte
ordinaire à un groupe à privilèges. Des outils publics de cartographie
calculent ces chemins~; les défenseurs les utilisent pour les couper.

\begin{tableaucours}{L{1.0cm}L{5.6cm}L{5.6cm}}{%
    \tableauHeader{Arête} &
    \tableauHeader{Relation} &
    \tableauHeader{Exemple chez LiZbiZ}
}
1 & Membre de & svc\_sql\_facturation \(\rightarrow\) Admins du domaine \\
2 & Session ouverte sur & Un administrateur connecté sur PC-GAR-12 \\
3 & Administrateur local de & Le groupe « Support Garoua » sur tous les
postes de Garoua \\
4 & Droit d'écriture sur & Le groupe « Support Garoua » peut modifier le
groupe « Admins serveurs » \\
\end{tableaucours}
\vspace{-0.8em}
\begin{center}{\small\itshape Tableau~10.3. Des relations qui forment des chemins}\end{center}

\begin{boxalert}
\textbf{Violation typique.} Corriger le seul compte exploité et déclarer le
problème résolu.

\smallskip
\textbf{Règle~:} le test a montré \emph{un} chemin. La correction consiste à
cartographier tous les chemins vers les groupes à privilèges et à couper les
plus courts, puis à refaire la cartographie.
\end{boxalert}

% ============================================================================
\section{Le Modèle par Niveaux et les Défenses}
\label{sec:c10-defense}
% ============================================================================

Le \termecle{modèle d'administration par niveaux} sépare les comptes et les
machines en niveaux (contrôleurs de domaine et identités~; serveurs~;
postes)~: un compte d'un niveau ne s'ouvre jamais sur une machine d'un
niveau inférieur, pour que ses secrets n'y soient jamais exposés.

\begin{tableaucours}{L{4.6cm}L{7.6cm}}{%
    \tableauHeader{Défense} &
    \tableauHeader{Effet}
}
Comptes de service gérés & Mots de passe longs, aléatoires, changés
automatiquement~: le kerberoasting devient inutile \\
Mots de passe de 25~caractères au moins pour les autres comptes de service
& Rend le cassage hors ligne irréaliste \\
Retrait des droits excessifs & Un compte de service n'est pas
administrateur du domaine \\
Pré-authentification obligatoire & Supprime l'AS-REP roasting \\
Mot de passe administrateur local unique (LAPS) & Bloque le mouvement
latéral (cours SEC3266) \\
Groupe Protected Users, modèle par niveaux & Limite l'exposition des
secrets à privilèges \\
Détection & Événements de demande de ticket inhabituels (nombreux SPN en
peu de temps, chiffrement faible demandé) \\
\end{tableaucours}
\vspace{-0.8em}
\begin{center}{\small\itshape Tableau~10.4. Les défenses de l'annuaire}\end{center}

\begin{boxobjet}{T-01 · Compte de service à privilèges et mot de passe faible}
\textbf{Description~:} « svc\_sql\_facturation » a un SPN, un mot de passe de
11~caractères prévisible, jamais changé, et appartient aux Admins du
domaine. \textbf{Preuve~:} appartenance au groupe (capture du 28~octobre,
14~h~08), mot de passe non reproduit. \textbf{Gravité~:} critique.
\textbf{Impact~:} contrôle total du domaine par tout utilisateur
authentifié. \textbf{Recommandation~:} retirer le compte des Admins du
domaine, le convertir en compte de service géré, cartographier les autres
chemins. \textbf{Détection~:} alerte sur les demandes de tickets en rafale.
\end{boxobjet}

\begin{boxmetier}{Analyste SOC et équipe rouge}
\textbf{Placer un compte « appât ».} Un compte de service fictif, avec un
SPN et aucun usage légitime~: toute demande de ticket le concernant signale
une tentative de kerberoasting.
\end{boxmetier}

\separateur

\begin{boxresume}
\textbf{Ce qu'il faut retenir de ce chapitre~:}
\begin{itemize}
    \item \textbf{L'annuaire} est lisible par tout utilisateur~; les
          groupes à privilèges, les comptes de service et les droits sont
          les cibles.
    \item \textbf{Kerberos} chiffre le ticket de service avec le secret du
          compte de service~: d'où le kerberoasting.
    \item \textbf{AS-REP roasting}~: la pré-authentification désactivée
          expose le compte.
    \item \textbf{On raisonne en chemins}~; on cartographie et on coupe les
          plus courts.
    \item \textbf{Comptes de service gérés, moindre privilège, niveaux,
          détection}~: les défenses.
\end{itemize}
\end{boxresume}

\begin{boxinfo}
\textbf{Points de contrôle de l'élève.} Je sais~: nommer les groupes à
privilèges~; décrire les quatre échanges Kerberos~; expliquer le
kerberoasting et l'AS-REP roasting~; lire un chemin d'attaque~; proposer
les défenses et une détection.
\end{boxinfo}

\bigskip

\begin{boxexo}[title={\ding{45}\quad Exercice~10.1. Les groupes \niveaubase}]
Citez trois groupes ou types de comptes d'un domaine qui intéressent en
priorité l'attaquant, et pourquoi.
\end{boxexo}

\begin{boxexo}[title={\ding{45}\quad Exercice~10.2. Le secret \niveaubase}]
Avec quel secret le ticket de service est-il chiffré~? Pourquoi est-ce la
source du kerberoasting~?
\end{boxexo}

\begin{boxexo}[title={\ding{45}\quad Exercice~10.3. Chronologie \niveaubase}]
Combien de temps s'écoule entre l'interrogation de l'annuaire et la preuve
d'accès aux Admins du domaine~?
\end{boxexo}

\begin{boxexo}[title={\ding{45}\quad Exercice~10.4. Les trois faiblesses \niveaumoyen}]
Identifiez les trois faiblesses de configuration combinées dans T-01, et la
mesure qui supprime chacune.
\end{boxexo}

\begin{boxexo}[title={\ding{45}\quad Exercice~10.5. Pourquoi le SOC n'a rien vu \niveaumoyen}]
Pourquoi le kerberoasting est-il discret~? Quelle règle aurait pu le voir~?
\end{boxexo}

\begin{boxexo}[title={\ding{45}\quad Exercice~10.6. Le chemin \niveaumoyen}]
À partir du Tableau~10.3, décrivez un second chemin, sans le compte de la
facturation, qui mène un membre de « Support Garoua » vers des droits
élevés.
\end{boxexo}

\begin{boxexo}[title={\ding{45}\quad Exercice~10.7. Niveaux \niveaumoyen}]
Un administrateur du domaine lit ses courriels sur PC-GAR-12. Quelle règle
du modèle par niveaux viole-t-il~? Quel risque crée-t-il~?
\end{boxexo}

\begin{boxexo}[title={\ding{45}\quad Exercice~10.8. AS-REP \niveaumoyen}]
Comment recenser les comptes dont la pré-authentification est désactivée~?
Que faire de chacun~?
\end{boxexo}

\begin{boxexo}[title={\ding{45}\quad Exercice~10.9. Le plan \niveauexpert}]
Rédigez le plan de durcissement de l'annuaire de LiZbiZ sur trois mois, à
partir du Tableau~10.4.
\end{boxexo}

\begin{boxexo}[title={\ding{45}\quad Exercice~10.10. Mesurer les chemins \niveauexpert}]
Proposez un indicateur fondé sur les chemins d'attaque (nombre, longueur)
et un objectif pour le comité des risques.
\end{boxexo}

\bigskip

\begin{boxinfo}
\textbf{Pour aller plus loin.}
\begin{itemize}
    \item RFC~4120~\cite{rfc4120}~; MITRE ATT\&CK~\cite{mitreattack}.
\end{itemize}
\end{boxinfo}

\begin{boxinfo}
\textbf{Transition.} Sur une machine, l'attaquant cherche aussi à passer
d'un compte limité à un compte d'administration~: le
Chapitre~\ref{chap:privileges} traite de l'élévation de privilèges.
\end{boxinfo}
